% ==================== 第1章 ====================
\chapter{绪论}
\label{ch:intro}

\section{研究背景与问题提出}

金融文本分析连接财务报告理解、事实检索与业绩判断，但真实查询并不服从单一的问答契约。确定信息问题要求系统从截止时点以前的10-K、10-Q、8-K及其表格、脚注和XBRL事实中读取或计算已经陈述的结果，评价重点是答案准确性与证据可定位性；前瞻性断言则涉及尚未实现的营收指引、产能投放、监管审批或合同履约，其真值在截止时点之后才会揭示，合理输出可能是点判断、可信集或弃权。管理层预测研究也表明，指引同时包含区间、时点和可验证性，不能等同于已实现业绩\citep{hirst2008management}。围绕长文档、表格推理和证据获取，金融问答与检索增强生成已经形成系统方法\citep{lewis2020rag,chen2021finqa,zhu2021tatqa,islam2023financebench,reddy2024docfinqa}，主动检索、自适应检索和智能体框架进一步优化了检索时机与证据使用\citep{jiang2023flare,jeong2024adaptiverag,yao2023react,shinn2023reflexion}；金融断言验证研究将陈述映射为支持、反驳或证据不足\citep{zhao2024findver,rangapur2025finfact}，选择性预测则通过拒识控制低置信度输出\citep{chow1970optimum,geifman2017selective}。这些研究分别改进了证据获取、断言核验和风险控制，却通常预先固定回答接口。固定契约用于混合负载时，确定问题可能被压缩为分类裁决，前瞻性断言也可能被阅读器输出为无保留的点答案。尚未被充分处理的关键问题因此不是单一检索器或分类器的性能，而是系统能否先判断一条查询应遵循事实检索契约还是前瞻性判断契约，并在相应契约下维护证据状态。本文据此提出FinPlanRAG：给定查询$q$与截止时点$c$，LLM路由器选择义务检索或AnchorEvolve；前者以类型化义务驱动时点取证并返回带证据编号的闭合答案，后者以稳定关系、行业模式和实例三层锚点组织前瞻性证据，依据可靠性、作用域和时效更新结果信念。点答案、可信集或弃权由校准读出生成，使契约选择、证据演化与风险控制在同一时点约束下形成可审计链条。

\section{研究内容与方法}

本文按照查询契约识别、分支证据建模和统一实证检验展开。首先，将金融混合查询形式化为确定信息问题与前瞻性断言问题的并集，并以截止时点界定可用信息，所有检索、特征构造、校准和读出选择均遵守时点约束，以控制前视偏差\citep{lopezdeprado2018advances}。其次，构建由LLM路由器、义务检索和AnchorEvolve组成的总体框架。路由器仅依据查询与截止时点元数据返回契约标签、置信度和理由；确定信息问题被解析为类型化义务，系统围绕未覆盖义务迭代检索，由LLM阅读器判断证据闭包并生成带编号答案；前瞻性断言被映射到有限结果空间，AnchorEvolve抽取L0稳定关系、L1行业模式和L2实例锚点，通过五类资格判定更新信念，校准读出据此生成点答案、可信集或弃权。最后，本文在DocFinQA、XBRL核验收入指引和FinDVer上，使用DeepSeek Flash、GLM 5.3 Flash和Qwen3-8B进行同基座比较，并分别检验路由效用、答案质量、证据成本、锚点增量信息、概率校准和非对称决策损失。这样的研究设计使契约识别、证据组织与风险控制能够被分层验证，同时保留端到端系统在混合负载上的整体评价。

\section{主要贡献}

本文的主要贡献形成于查询层、证据层和评价层的连续设计。查询层将回答形式选择表述为混合负载上的契约路由，把“事实在截止时点是否已经存在”置于检索与生成之前，使闭合问答与前瞻性裁决可以共享时点约束而保持不同输出接口。证据层提出AnchorEvolve的金融断言实例化，以L0稳定关系、L1行业模式和L2实例锚点记录证据的取值、作用域、来源、时间与可靠性，通过五类资格判定形成可更新、可追溯的信念状态，并让接受的证据按类型和公开揭示日更新Beta可靠性；类型库扩展作为验证集门控的训练步骤，冻结后才用于测试。义务检索以类型化义务和阅读器充分性控制确定问题的证据获取。决策层用全动作损失感知读出把点答案、非单元素裁决集和弃权放在同一个期望损失下比较。评价层在DocFinQA、227条收入指引断言、470条A股年报断言和150条FinDVer断言上，以三套基座和同基座基线检验路由、义务检索、锚点推断、后验融合、类型库扩展、数值桥和外部迁移。由此，FinPlanRAG不仅给出端到端结果，也保留证据编号、锚点来源、资格判定和信念变化，为金融分析中的复核、问责与事后审计提供可复用接口。

% ==================== 第2章 ====================
\chapter{文献综述}
\label{ch:related}

本章围绕金融问答与断言验证、检索增强生成与自适应检索、选择性预测与可信集、锚点约束与时点证据四条线索梳理相关研究，并在此基础上归纳研究缺口与关键挑战。

\noindent\textbf{一、金融问答与断言验证}

金融问答基准测试模型在申报文件、表格和报告上的数值推理能力。FinQA研究基于财务报表的数值推理\citep{chen2021finqa}，TAT-QA研究表格与文本的混合推理\citep{zhu2021tatqa}，FinanceBench研究真实金融文档上的开放问答\citep{islam2023financebench}，FinBen研究金融领域的综合评测\citep{xie2024finben}。DocFinQA在FinQA问题的基础上加入完整10-K上下文，把证据选择从段落级扩展到长文档级\citep{reddy2024docfinqa}。这些基准的共同点是答案在文档中已经存在，评价指标主要是精确匹配或词级F1。它们衡量的是阅读器能否在给定文档中找到并计算答案，而不是系统能否判断查询是否应当以闭合答案回应。

金融申报文件的结构与通用问答语料不同。10-K和10-Q通常包含多级标题、跨页表格、脚注和附注，同一指标可能在正文、表格和管理层讨论中以不同口径出现；8-K则按事项编号披露业绩、指引、并购和监管事件。XBRL数据提供结构化财务事实，但季度口径、单位、重述和可用日期仍需与文本对齐。通用阅读器擅长语义匹配，却不一定能处理期间、单位和证据闭合，因此金融问答需要把文档结构、时点约束和数值计算纳入同一证据策略。

断言验证基准测试模型判断一条陈述是否被证据支持、反驳或无法判断。FinDVer面向金融领域的断言验证，要求模型在长文档和混合内容中给出裁决和解释\citep{zhao2024findver}；Fin-Fact面向证据驱动的金融事实核查，强调多模态证据和解释生成\citep{rangapur2025finfact}。这类任务的输出是离散裁决，评价指标是裁决准确率和解释质量。验证器通常假设陈述已经被给定，且证据已经被检索，因而没有回答“这条陈述是否应当进入验证流程”的问题。

近期金融RAG系统进一步加入类型化证据、期间约束、来源追踪和语料对齐检索\citep{mansouri2026finrank,jiang2026finrate,chen2026hcrag,chi2026finsagent,zhou2026fincards}。这些系统在证据表示和检索配方上明显优于通用RAG，但输出接口仍通常是固定答案或固定裁决。它们改善了到达阅读器的证据，却没有改变阅读器与验证器之间的契约边界。对于同时包含事实问题和前瞻性断言的金融分析场景，这种固定接口仍然是主要限制。

\noindent\textbf{二、检索增强生成与自适应检索}

RAG将检索器与生成器组合，假设答案可以由检索到的上下文拼装\citep{lewis2020rag}。主动检索和自适应检索进一步决定何时检索或检索多少上下文\citep{jiang2023flare,jeong2024adaptiverag}，智能体框架则交错进行推理、检索和自我修正\citep{yao2023react,shinn2023reflexion}。Self-RAG通过自我反思决定是否检索和如何评价证据\citep{asai2024selfrag}，CRAG在检索质量不足时触发纠正性检索\citep{yan2024crag}，FLARE在生成过程中按置信度触发检索\citep{jiang2023flare}，Adaptive-RAG按问题复杂度选择检索策略\citep{jeong2024adaptiverag}。

这些方法改善了检索过程和证据使用，但输出通常仍是生成的答案。它们可以决定“是否检索”和“检索多少”，却不能决定“返回答案、可信集还是弃权”。在确定信息问题上，这种输出是合适的；在前瞻性断言上，单一生成答案无法表达结果分布、证据作用域和决策风险。检索策略的自适应并不等于输出契约的自适应，二者是两个不同层次的决策。

本文的义务检索分支借用了自适应检索的迭代思想，但把停止条件从“检索轮数”改为“义务覆盖与阅读器充分性”。这一改动使证据预算服务于问题的语义结构，而不是服务于固定的$k$。AnchorEvolve则把输出契约从生成答案扩展为点答案、可信集和弃权的组合，使系统可以在证据不足时保留不确定性。

\noindent\textbf{三、选择性预测、校准与可信集}

拒识选项在识别错误和拒识率之间进行权衡\citep{chow1970optimum}，选择性分类通过置信度阈值在低置信度输入上弃权\citep{geifman2017selective}。现代分类器通常存在校准误差，温度缩放是常用基线\citep{guo2017calibration}；代价敏感学习为不同错误赋予不同价格\citep{elkan2001cost}；分裂共形预测则把得分转换为在可交换性下具有分布无关覆盖率的集合\citep{vovk2005algorithmic,angelopoulos2021gentle}。分类共形方法进一步构造自适应集合\citep{romano2020classification}，分布漂移下的自适应共形方法则允许覆盖率随时间调整\citep{gibbs2021adaptive}。

FinAbstain进一步把时点检索、多模态证据、证据核验、不确定性校准和显式弃权组合成选择性金融预测流程\citep{torres2026finabstain}。它是与本文模糊分支最接近的工作，但仍然把查询预先视为预测任务，没有在闭合问答与前瞻性裁决之间做契约路由；其证据核验也是不确定性聚合之前的过滤，而不是独立的多类别资格判定阶段。

这些方法提供了损失、置信度和集合构造的工具，但通常假设任务和回答接口已经固定。选择性分类器在“答案”和“弃权”之间选择，验证器在“支持”“反驳”和“证据不足”之间选择，共形方法在给定标签空间上构造集合。它们没有解决“标签空间本身是否适用于当前查询”的问题。对于前瞻性断言，标签空间是结果桶，证据是锚点，输出还需要说明证据的作用域和时效；这些结构不在标准选择性预测的范围内。

本文借用选择性预测的损失框架和共形集合的覆盖思想，但把标签空间、证据资格和输出动作放在同一个锚点状态上。校准随机森林和温度缩放用于控制后验概率，损失感知门用于在答案、可信集和弃权之间选择。与标准共形方法相比，本文的可信集由后验质量构造，覆盖率在实验中被报告，但并未假设交换性；这是金融时点数据与标准共形设定之间的差异。

\noindent\textbf{四、锚点、部分识别与时点证据}

部分识别表示与现有证据一致的结果集合\citep{manski2003partial}。在金融机器学习中，模型决策只能使用当时可得的信息，否则会产生前视偏差\citep{lopezdeprado2018advances}。本文借用约束和可行集的思想，将锚点定义为来源与假设之间的关系：硬锚点可以排除某个结果桶，软锚点只重新加权结果信念；证据过滤和模式过滤均按截止时点执行。

与部分识别相比，本文的锚点不是完整的结构模型，而是由LLM从文本中抽取的语义约束。锚点带有类型、取值、作用域、来源、时间、可靠性和有效期，因而可以被逐条检查。部分识别给出的是在给定模型下与证据一致的参数集合，锚点给出的是在给定断言下与证据一致的结果桶约束。两者都强调“证据不足时不要强行点识别”，但本文更关注证据的可审计性和可更新性。

时点约束贯穿本文的两个分支。义务检索只接受可用日期不晚于截止时点的文档，AnchorEvolve只接受截止时点之前的证据，随机森林只使用截止时点之前可得的特征，校准和权重选择只在开发集与验证集上进行。这一约束使“当时可得信息”成为系统设计的一部分，而不是事后评价的附加条件。

\noindent\textbf{五、研究缺口与关键挑战}

现有研究分别推进了长文档检索、金融问答、断言验证和低置信度拒识，但尚未把查询契约选择与前瞻性证据演化纳入同一框架。本文认为，这一缺口可以分解为三个相互衔接的挑战。三个挑战分别对应查询层、证据层和决策层，也构成本文框架的设计依据。

第一个挑战是查询契约的识别。金融查询的表面形式与真实契约并不一一对应。“预期股利”“未来支付金额”等表述包含未来时态，但答案可能在文件中已经给出；管理层指引包含精确数字，却是关于尚未实现结果的前瞻性断言。系统如果按关键词或时态路由，会把确定问题送入裁决分支，或把前瞻性断言送入阅读器。真正的判断依据是查询所要求的事实是否在截止时点已经存在，以及输出应当采用闭合答案还是可校准判断。这一挑战的难点不在分类器本身，而在于契约是查询、证据与时点之间的关系，而不是查询文本的主题标签。

第二个挑战是前瞻性证据的状态化。现有验证器通常把断言直接映射为支持、反驳或证据不足，阅读器则直接生成答案。两者都没有保留证据的作用域、可靠性、时效和来源，也没有在证据变化时更新判断。对于前瞻性断言，证据往往只支持某个结果桶、只在某个时间范围内有效、只对某个公司或行业成立；把这些信息压缩为一个标签，系统就无法解释判断为何改变，也无法在结果揭示后更新可靠性。本文认为，前瞻性断言需要一个可审计的中间状态，使证据到信念、信念到决策的每一步都可以追溯。

第三个挑战是语义判断与风险校准的分工。LLM能够从文本中抽取锚点、解释证据和形成语义判断，但其概率输出往往存在类别偏置和校准误差；统计模型可以使用结构化特征进行概率校准，却缺少对文本证据的理解。若让LLM直接输出最终裁决，系统可能过度自信；若只用统计模型，系统又无法利用文本中的语义约束。本文认为，合理的分工是让LLM构造语义状态，让校准程序控制决策风险，并在验证集上选择读出方式，而不是在测试集上比较后择优。

上述三个挑战对应本文的三项贡献。契约识别由查询契约路由回应；证据状态化由AnchorEvolve的三层锚点、五类资格判定和信念更新，以及义务检索回应；语义与校准的分工由校准后端、损失感知门和可信集读出回应。三项贡献分别对应查询层、证据层和决策层，构成逐层回应。

% ==================== 第3章 ====================
\chapter{研究设计：研究问题、假设与总体框架}
\label{ch:framework}

\section{研究问题}

金融文本分析中的查询可以分为两类。确定信息问题要求返回申报文件在截止时点已经陈述或可以计算的事实，答案通常是数值、段落或明确的预测。前瞻性断言问题要求判断尚未实现的未来结果，其真值取决于截止时点之后披露的信息，正确输出可能是一个点答案、一个可信集，或者在证据不足时弃权。设查询为$q$，截止时点为$c$，工作负载是确定问题集合$\mathcal Q_D$与前瞻性断言集合$\mathcal Q_V$的混合。系统调用LLM路由器
\begin{equation}
r(q,c)\in\{D,V\}
\end{equation}
判断查询属于哪一类，再将确定查询送入义务检索分支，将前瞻性断言送入锚点演化分支。

两个分支返回不同的输出形式。义务检索分支返回闭合答案$a$、置信度与证据编号，证据不足时返回\texttt{insufficient}。AnchorEvolve返回点答案$h^*$、可信集$R_\alpha$、动作$\pi$与信念$b$，动作可以是回答、可信集裁决或弃权。路由器只读取原始查询和截止时点元数据，不读取“这是断言”的任务标签，也不读取测试结果。

上述问题可以展开为三个研究问题：
\begin{enumerate}[label=RQ\arabic*.,leftmargin=2em,itemsep=2pt,topsep=2pt]
\item 契约识别：仅依据查询与截止时点元数据，系统能否在混合负载上区分确定信息契约与前瞻性断言契约？
\item 确定信息取证：类型化义务与覆盖驱动的证据获取，能否在更低证据成本下提高确定信息问答的准确率与覆盖率？
\item 前瞻性裁决：三层锚点状态与校准读出，能否在保持分类准确率的同时降低决策损失，并提供可信集或弃权？
\end{enumerate}
三个问题分别对应查询层、证据层和决策层。RQ1检验路由器能否识别契约；RQ2检验义务检索能否以更低成本提高确定问答的准确率和覆盖率；RQ3检验锚点状态与校准读出能否把语义信息转换为低风险决策。

\section{研究假设}

三个研究问题对应三个可检验假设。假设的表述同时给出检验口径和否证条件，使后文实验能够直接检验框架是否有效，而不是只描述框架结构。

\textbf{H1（契约可识别性）}：仅依据查询与截止时点元数据，LLM路由器能够在混合负载上区分确定信息契约与前瞻性断言契约；在较强基座上，路由混合效用接近oracle，并高于Always-RAG和Always-Anchor。检验口径为573条混合负载上的路由准确率、混淆矩阵和混合效用。若路由器在多数基座上低于较好的固定分支，则H1不成立。

\textbf{H2（义务取证的效率）}：类型化义务与覆盖驱动的证据获取，能够在低于固定16段的平均证据成本下，提高DocFinQA的精确匹配和答案覆盖率，并迁移到FinDVer。检验口径为DocFinQA的300题固定样本和FinDVer testmini上的准确率、宏F1、平均段落数和覆盖率。若分支在多数基座上不高于固定16段阅读器，或证据成本没有下降，则H2不成立。

\textbf{H3（语义状态与校准读出的分工）}：三层锚点状态能够提供结构化语义信号，校准后端与损失感知门能够把该状态转换为低风险决策；最终系统在保持普通准确率的同时降低非对称决策损失，并在锚点消融和输出形式消融中表现出可解释的差异。检验口径为227条XBRL核验断言上的准确率、平衡准确率、宏F1、ECE、Brier、决策损失、可信集覆盖率和平均集合大小。若最终系统在所有基座上的决策损失不低于最强具名基线，或去除锚点后宏F1没有下降，则H3不成立。

三个假设共同构成对研究问题的可否证回答。H1对应查询层，H2对应确定信息的证据层，H3对应前瞻性断言的证据层与决策层。第\ref{ch:empirical}章按这三个假设组织实验，而不按数据集顺序堆叠结果。

\section{总体框架}

FinPlanRAG由三个模块组成。路由器是一个独立的LLM调用，负责判断查询契约；义务检索分支负责确定信息问题；AnchorEvolve分支负责前瞻性断言。路由器与两个分支之间通过结构化输出连接，路由器不读取分支输出，分支不改变路由结果。这样的结构使路由错误和分支错误可以被分开诊断：路由错误表现为查询被送入不匹配的契约，分支错误表现为在正确契约下的证据或判断错误。

路由器读取查询文本与截止时点元数据，返回路由标签、置信度和简短理由。提示词定义两个契约：确定信息请求要求返回申报文件在截止时点陈述或可计算的事实、数值、段落、估计或预测；前瞻性断言要求判断尚未实现的未来结果，例如实际营收是否落在指引区间内、产品是否按时发布、目标是否达成。提示词明确说明管理层指引是前瞻性断言，不是已实现事实，即使其中包含数字。

路由器不读取已实现结果，也不读取两个分支的输出。它只看到查询和截止时点元数据，因此在部署时可以按查询缓存，成本是一次短LLM调用。路由器的错误类型也具有语义含义：确定问题被送到AnchorEvolve，通常是因为查询文本包含未来时态或预期措辞；前瞻性断言被送到义务检索，通常是因为断言中包含可检索的数字或期间。这些错误为改进路由器提供了可解释的信号。

两个分支共享路由器，但不共享输出契约。义务检索分支返回闭合答案，AnchorEvolve返回点答案、可信集或弃权。共享路由器使系统只需一次契约判断；不共享输出契约使每个分支可以针对自己的证据结构和风险结构设计。图\ref{fig:framework}给出总体结构，第\ref{ch:determinate}章和第\ref{ch:anchor}章分别给出两个分支的完整设计。

\begin{figure}[H]
\centering
\resizebox{0.98\linewidth}{!}{%
\begin{tikzpicture}[
  font=\normalsize,
  >=Latex,
  box/.style={draw, rounded corners=3pt, align=center, text width=2.6cm, minimum height=0.9cm, inner sep=3pt, fill=white},
  panel/.style={draw, rounded corners=6pt, inner sep=7pt, fill=gray!5},
  arrow/.style={->, thick, draw=black!65},
  lbl/.style={font=\footnotesize, text=black!70}
]
\node[box, fill=gray!10] (q) at (0,0) {查询 $q$\\截止时点 $c$};
\node[box, fill=gray!18, text width=2.6cm] (r) at (2.7,0) {LLM 路由器\\契约判断};

\node[box] (o1) at (5.6,2.0) {类型化义务};
\node[box] (o2) at (8.6,2.0) {时点检索};
\node[box] (o3) at (11.6,2.0) {LLM 阅读器};

\node[box] (a1) at (5.6,-2.0) {三层锚点};
\node[box] (a2) at (8.6,-2.0) {资格与信念};
\node[box] (a3) at (11.6,-2.0) {校准与损失门};

\begin{scope}[on background layer]
\node[panel, fit=(o1)(o3), label={[font=\bfseries]above:义务检索分支}] {};
\node[panel, fit=(a1)(a3), label={[font=\bfseries]below:AnchorEvolve 分支}] {};
\end{scope}

\node[box, text width=2.1cm] (ans) at (14.6,2.0) {闭合答案};
\node[box, text width=2.1cm] (pt) at (14.6,0.6) {点答案};
\node[box, text width=2.1cm] (cs) at (14.6,-0.6) {可信集};
\node[box, text width=2.1cm] (ab) at (14.6,-2.0) {弃权};
\node[lbl, above] at (14.6,2.65) {输出};

\draw[arrow] (q) -- (r);
\draw[arrow] (r.east) -- ++(0.7,0) |- (o1.west);
\draw[arrow] (r.east) -- ++(0.7,0) |- (a1.west);
\draw[arrow] (o1) -- (o2);
\draw[arrow] (o2) -- (o3);
\draw[arrow] (a1) -- (a2);
\draw[arrow] (a2) -- (a3);
\draw[arrow] (o3) -- (ans);
\draw[arrow] (a3) -- (pt);
\draw[arrow] (a3) -- (cs);
\draw[arrow] (a3) -- (ab);

\draw[dashed, draw=black!45] (0.9,3.0) -- (0.9,-3.0);
\node[lbl, above] at (0.9,3.05) {时点约束 $V_c$};
\end{tikzpicture}
}
\caption{FinPlanRAG总体框架。路由器把查询送入义务检索分支或AnchorEvolve；前者返回闭合答案，后者返回点答案、可信集或弃权。左侧虚线表示时点约束。}
\label{fig:framework}
\end{figure}


\section{时点约束与评价协议}

系统在截止时点$c$形成判断，可用信息由可采纳证据集合$V_c$给出。设证据$e$的可用时间为$a(e)$，则
\begin{equation}
V_c=\{e:a(e)\le c,\ e\text{在}c\text{时未被撤回}\}.
\end{equation}
策略$\pi$在$c$时有效，当且仅当它是$\sigma(V_c)$可测的。这一条件排除了使用未来结果、未来披露或未来修订的证据。路由器不读取结果，锚点抽取器不接收结果，随机森林只使用截止时点之前可得的特征，校准和权重选择只在开发集与验证集上进行。测试结果不参与任何拟合或选择。

评价口径与输出契约对应。确定信息问题使用精确匹配、词级F1、答案覆盖率和证据成本；前瞻性断言使用准确率、平衡准确率、宏F1、MCC、AUROC、AUPRC、ECE、Brier、决策损失、可信集覆盖率和平均集合大小；混合负载使用路由准确率、混淆矩阵和混合效用。所有任务共享同一时点约束和同一套测试集不参与选择的原则，因此不同分支的结果可以在同一框架下比较。

评价协议由三部分组成。第一，主结果使用配对自助法给出置信区间，准确率差异使用McNemar检验；所有区间在断言或问题级别重采样。第二，随机森林使用五个折种子，测试年份的训练集来自其他公司折，且结果日期早于测试年份1月1日；测试公司不出现在训练集中。第三，所有LLM响应按请求负载的哈希缓存，提示词、运行脚本、逐条预测和表格生成脚本随论文发布，API密钥只从环境变量读取。

这一协议的目的是让每个假设都有对应的证据边界。H1对应573条混合负载；H2对应DocFinQA的300题固定样本和FinDVer testmini；H3对应227条XBRL核验断言、锚点消融和输出形式消融。

% ==================== 第4章 ====================
\chapter{义务检索分支：Obligation Retrieval}
\label{ch:determinate}

义务检索分支可以概括为四步：把问题表示为类型化义务，构建时点证据空间，按未覆盖义务迭代取证，由阅读器生成带证据编号的闭合答案。四步依次对应本章四节。

\section{状态表示：类型化义务}

义务检索分支把自然语言问题绑定为可检索、可核对的类型化义务。义务元组定义为
\begin{equation}
o=(e,f,y,p,m,u),
\end{equation}
其中$e$为实体，$f$为表单或文档类型，$y$为财年，$p$为财期，$m$为指标，$u$为单位。以“某公司2023财年第四季度营业收入是多少”为例，实体是公司，表单是10-K或10-Q，财年是2023，财期是第四季度，指标是营业收入，单位是美元。义务项从问题和查询元数据中抽取；抽取器只负责识别约束，不负责回答。

文档只有在可用日期不晚于截止时点、且元数据与义务一致时才被接受。元数据一致性检查覆盖实体、表单、财年、财期、指标和单位。与只按问题文本检索相比，义务元组把实体、期间和单位变成可检查的约束。例如，问题文本中的“2023年”可能指财年、自然年或季度，BM25无法区分；义务元组把“2023财年第四季度”固定为可核对的期间，从而减少跨期文档和单位不一致带来的错误。

类型化义务还提供了证据覆盖的度量。设问题$q$的义务项集合为$T(q)$，已获取证据集合为$A_t$，则未覆盖义务项为
\begin{equation}
U_t=T(q)\setminus\{t\in T(q):\exists a\in A_t,\ t\in a\}.
\end{equation}
未覆盖项既用于构造下一轮检索查询，也用于判断何时停止。义务覆盖不是答案正确性的充分条件，但它把“证据是否覆盖了问题的关键约束”变成可观察的状态，使系统可以在证据不足时继续检索或返回证据不足，而不是在固定轮数后强行生成答案。

\section{证据空间：时点语料与块索引}

分支把文档切分为段落块和表格块。段落块按行合并，单块不超过1200字符；表格块按连续表格行合并，单块不超过4000字符。每个块记录文档编号、可用时间、文本、词项和元数据。切块的目标是保留表格行与表头的对应关系，同时避免单块过长导致阅读器输入超出预算。对于跨页表格，分支保留表头并在块内重复列名，以减少阅读器对列含义的误判。

分支使用BM25建立索引，按问题与未覆盖义务项检索候选块。BM25的优势是无需训练、可复现、对长文档稳定；缺点是不理解语义，因此需要义务覆盖和阅读器反馈补足。检索查询由原问题与$U_t$拼接而成，候选块按得分排序，去重后最多新增两个块。这样做的原因是：一次加入过多块会稀释阅读器的注意力，也会增加token成本；一次加入过少块则可能需要更多轮检索。两块的设置是在证据覆盖和调用成本之间的折中。

时点过滤在检索前执行。只有可用日期不晚于截止时点的文档进入索引；被后续版本撤回或替代的文档被排除。对于同一文档的多个版本，分支保留截止时点之前的最新版本。时点过滤使系统不会使用未来披露的信息，也避免了“文档存在但当时尚未发布”的前视偏差。

\section{取证策略：覆盖驱动的迭代检索}

第$t$轮检索后，分支把新增块加入证据集$A_t$，调用LLM阅读器判断义务是否闭合。停止条件有三类：阅读器报告证据充分且$U_t=\varnothing$；阅读器给出非空答案、置信度不低于0.8且$U_t=\varnothing$；预算$B$耗尽。默认$B=8$，并报告$B=12$和$B=16$。前两类条件把停止权交给证据状态和阅读器判断，第三类条件提供成本上限。

覆盖驱动的取证把停止变量从检索前的固定预算移到义务覆盖和阅读器充分性上。对于简单问题，分支可以在两三轮内闭合；对于复杂问题，分支可以继续检索直到预算耗尽。这种自适应性使平均证据成本低于固定预算，同时在难题上保留足够的证据预算。

置信度门控是覆盖驱动取证的补充。当阅读器置信度低于0.8或报告证据不足时，分支回退到固定16段证据集重新阅读。回退不是对每个问题支付16段成本，而是对规划器无法闭合的问题支付更高成本。实验表明，回退只在规划器无法闭合的问题上触发，DeepSeek Flash上约为24\%，GLM 5.3 Flash上比例更高；它弥补了规划器单独的准确率缺口，同时保持平均段落数低于固定16段。完整流程见算法\ref{alg:obligation}。



\section{答案生成：阅读器与闭包}

阅读器接收问题与已获取块，返回结构化答案
\begin{equation}
(\text{answer},\text{confidence},\text{sufficient},\text{missing},\text{evidence\_ids},\text{rationale}).
\end{equation}
答案要求为短字符串；若证据不足，返回\texttt{insufficient}并将\texttt{sufficient}置为\texttt{false}。分支把阅读器的\texttt{missing}字段拼入下一轮查询，形成阅读器驱动的取证。例如，阅读器报告缺少“第四季度”或“单位”，下一轮查询会加入这些词项，从而把检索引向缺失的证据。

闭包规则要求答案附带证据编号。每个答案都可以追溯到具体段落或表格块，证据不足时有明确的缺失信息。这一规则使义务检索分支的输出可以被下游审计：读者可以检查答案是否由证据支持，也可以检查系统在证据不足时是否选择弃权而不是猜测。

解析失败在评价中计为错误。阅读器被要求返回JSON，但基座可能返回解释性文字或格式错误的JSON。分支对解析失败不进行静默修复，而是记录为错误，以避免评价口径被解析器放宽。Qwen3-8B在DocFinQA上返回较多解释性文字，导致精确匹配接近地板；这一结果作为基座能力边界报告，不用于义务检索分支的主比较。

\begin{algorithm}[H]
\caption{义务检索的单轮取证}
\label{alg:obligation}
\begin{algorithmic}[1]
\Require 问题$q$，截止时点$c$，义务项$T(q)$，预算$B$
\Ensure 答案$a$或\texttt{insufficient}
\State $A\gets\varnothing$; $U\gets T(q)$
\For{$t=1,\dots,B$}
  \State $C_t\gets\textsc{BM25}(q\cup U,\,V_c)$
  \State $A\gets A\cup\textsc{Top2}(C_t\setminus A)$
  \State $(a,\text{conf},\text{suff},\text{miss},\text{ids})\gets\textsc{Reader}(q,A)$
  \State $U\gets T(q)\setminus\textsc{Covered}(A)$
  \If{$\text{suff}=\text{true}$ and $U=\varnothing$}
    \State \Return $a$
  \EndIf
  \If{$\text{conf}\ge0.8$ and $a\ne\varnothing$ and $U=\varnothing$}
    \State \Return $a$
  \EndIf
  \State $q\gets q\cup\text{miss}$
\EndFor
\State $A\gets\textsc{Fixed16}(q,V_c)$
\State $(a,\text{conf},\text{suff},\dots)\gets\textsc{Reader}(q,A)$
\If{$\text{suff}=\text{true}$}
  \State \Return $a$
\Else
  \State \Return \texttt{insufficient}
\EndIf
\end{algorithmic}
\end{algorithm}

义务检索的时间成本主要由BM25检索和阅读器调用组成。BM25在本地索引上完成，不产生API成本；阅读器调用次数由停止条件决定，最坏情况下为预算$B$加上一次固定16段回退。设文档块数为$N$，单次BM25检索的复杂度为$O(N\log N)$，相对于LLM调用可以忽略。

可审计性来自证据编号和缺失信息。每个答案都附带实际使用的段落或表格块编号，证据不足时返回具体的缺失项。审计者可以检查答案是否由证据支持，也可以检查系统在证据不足时是否选择弃权而不是猜测。

% ==================== 第5章 ====================
\chapter{AnchorEvolve：前瞻性断言分支}
\label{ch:anchor}

AnchorEvolve是FinPlanRAG的前瞻性断言分支。本章先给出状态表示、证据来源和总体结构，再给出完整算法，包括证据构造、锚点抽取、五类资格判定、时间演化、类型库扩展、信念更新、全动作读出、数值桥和披露日期断言。

% ==================== 第4章 ====================

\section{设计原则}

FinPlanRAG的设计遵循六条原则，其核心模块AnchorEvolve承担证据状态化、资格判定、时间演化和信念更新。这些原则不是附加的功能列表，而是对第\ref{ch:framework}章所述研究缺口的直接回应。

第一，契约先行。系统在检索证据之前先确定公司、目标期间、指标、结果空间和截止时点。契约把“回答什么”写成显式对象，使不同方法可以在同一口径下比较，也使用户可以在运行前检查阈值是否合理。契约不是系统的内部细节，而是审计报告的第一部分。

第二，来源落地。每条锚点必须引用一个具体的证据来源，并给出逐字引文。系统不接受没有来源的判断，也不接受只在理由中描述证据而不给出引文的锚点。来源落地使锚点可以被逐条复核，也使错误的锚点可以被定位到具体的网页或段落。

第三，资格先于权重。锚点在进入信念更新之前先经过资格判定。空泛、过期、对抗三类锚点不进入更新；硬锚点和软锚点分别处理。这个顺序很重要：如果先把所有锚点加权再过滤，被拒绝的锚点仍然会影响后验；如果先过滤再加权，每个被采纳的锚点都有明确的资格理由。

第四，揭示后演化。锚点类型的可靠性只在结果公开揭示之后更新，且只使用揭示日早于当前截止时点的结果。这个约束使可靠性演化在回测和部署中具有相同的含义，避免把未来结果泄漏到时点判断。

第五，损失感知。输出不是固定的点答案，而是在点答案、非单元素裁决集和弃权三个动作上最小化期望损失。损失函数把过度乐观、漏报、集合宽度和弃权成本写成显式参数，使不同风险偏好可以在同一框架下表达。

第六，审计优先。系统把证据、锚点、资格、可靠性、信念轨迹、动作损失和日期估计写入结构化文件。审计不是事后附加的日志，而是系统输出的一部分。没有审计文件的判断不能被本文视为可复核的判断。

\section{三层锚点状态}

FinPlanRAG的核心状态是AnchorEvolve的三层锚点。三层分别是稳定关系层（L0）、行业模式层（L1）和实例层（L2）。这里的锚点不是知识图谱中的实体或关系，也不是行为金融中的anchoring bias；L0、L1和L2分别对应关系约束、行业模式绑定和带来源的证据实例，锚定的是证据对结果的约束关系，而不是一个实体节点。AnE用“anchor”指训练数据的真值锚定，本文的锚点指证据状态中的结构化对象，二者机制不同\citep{wang2026ane}。

L0是稳定关系本体。它定义一组可复用的关系类型，包括界、先后、前置时间、持续时间、守恒、比率、单位换算和作用域。L0不随行业变化，为锚点提供语义骨架。例如，“当前季度值落在指引区间内”是一个界关系，“新产品在指引季度之后投产”是一个先后关系，“产能爬坡需要两个季度”是一个前置时间关系，“分部收入之和等于总收入”是一个守恒关系，“毛利率等于毛利润除以营业收入”是一个比率关系。L0的作用是限制锚点可以表达的关系种类，使不同行业的锚点共享同一套语义。

L1是行业模式。每个锚点类型是一个四元组
\begin{equation}
\tau=(\text{名称},\ \text{L0关系},\ \text{字段},\ \text{生效轮次}),
\label{eq:l1}
\end{equation}
其中字段给出该类型需要从证据中读取的信息，生效轮次给出该类型可以在哪一轮断言中使用。收入指引的初始模式包括指引宽度、当前季度增长、公司历史、行业历史、文本信号、产能与需求、分析师预测、披露历史和监管规则。L1把L0关系绑定到具体字段和生效规则。例如，“指引宽度”绑定到界关系，字段是上限与下限之差除以中点；“当前季度增长”绑定到比率关系，字段是当前季度值相对指引中点的比例；“披露历史”绑定到先后关系，字段是公司过去若干年年报的公告日期。

L2是锚点实例。每条锚点是一个元组
\begin{equation}
a=(id,\tau,v,u,s,src,d,r,\text{hard},\hat h,\text{note}),
\label{eq:l2}
\end{equation}
其中$id$是锚点编号，$\tau$是L1类型，$v$是取值或区间，$u$是单位，$s$是作用域，$src$是来源编号，$d$是来源日期，$r$是可靠性，$\text{hard}$是硬锚点标记，$\hat h$是隐含结果标签，$\text{note}$是理由。L2是LLM抽取器的直接输出，也是资格判定和信念更新的输入。每条L2锚点都记录来源和时间，因此可以被逐条检查。

三层的关系可以概括为：L0定义可以说哪些关系，L1定义某类断言需要哪些字段，L2记录某条证据的具体取值。三层不是三个独立训练的模块，而是同一状态的不同抽象层次。L1类型用于可靠性演化，L0关系用于语义分类，L2实例进入信念更新。这样的分层使系统可以在不改变L0的情况下增加新的L1类型，也可以在不改变L1的情况下替换L2抽取器。

\section{锚点类型}

表\ref{tab:anchortypes}给出本文使用的主要锚点类型。类型的选择遵循两个标准：一是该类型在金融文本中有稳定的语言表达，二是该类型能够对应到结果空间中的一个或多个标签。类型不是越多越好；过多的类型会增加抽取器的负担，也会增加可靠性后验的稀疏性。

\begin{table}[H]
\centering
\zihao{5}
\caption{FinPlanRAG的主要锚点类型。}
\label{tab:anchortypes}
\begin{tabular}{p{2.6cm}p{2.4cm}p{7.2cm}}
\toprule
类型 & L0关系 & 作用 \\
\midrule
指引宽度 & 界 & 指引区间越宽，结果落在区间内的先验越高；区间本身不是结果落在区间内的证据。\\
当前季度增长 & 比率 & 当前季度实际值相对指引中点的比例，用于判断全年结果的起点。\\
公司历史 & 先后 & 公司过去已揭示结果相对指引的偏差，用于估计公司层面的兑现倾向。\\
行业历史 & 先后 & 同行业其他公司已揭示结果的分布，用于估计行业层面的共同冲击。\\
文本信号 & 语义 & 管理层评论、需求描述、订单和产能信息，用于补充结构化字段。\\
产能与需求 & 前置时间 & 产能投放、利用率和订单覆盖，用于判断收入增长的可持续性。\\
分析师预测 & 界 & 卖方机构对目标期间的预测，作为软锚点而不是事实。\\
披露历史 & 先后 & 公司过去年度报告的公告日期，用于披露日期断言。\\
监管规则 & 界 & 交易所的披露截止日，作为上界而不是具体日期证据。\\
\bottomrule
\end{tabular}
\end{table}

锚点类型与结果标签的关系不是一一对应。一条“当前季度增长”锚点可能支持“高于”标签，也可能只支持“不低于”某个下界；一条“指引宽度”锚点通常不直接支持任何标签，而是影响区间内的先验。资格判定和信念更新需要分别处理这些情况。锚点的$\hat h$字段记录抽取器认为该锚点隐含的标签；如果锚点不提供标签信息，$\hat h$为空，锚点被判为空泛。

\section{证据来源与作用域}

FinPlanRAG使用三类证据来源。

第一类是公司申报文件和公告。美国公司使用8-K、10-Q、10-K和XBRL数据；中国A股公司使用巨潮资讯网的公告记录、年报、季报和临时公告。申报文件提供结构化事实和原文段落，公告记录提供披露日期和文件链接。官方来源的优先级高于媒体转载，因为官方来源的日期和文本更可靠。

第二类是财经媒体和研究报告。媒体文章可以提供管理层评论、行业需求、产能和价格信息，但可能存在转载、改编和日期错误。系统要求媒体来源给出逐字引文，并对无日期来源降权。研究报告中分析师预测作为软锚点，不作为事实。

第三类是交易所规则和监管文件。这类来源提供披露截止日、披露格式和合规要求。监管规则通常只给出上界，不能直接支持某个具体日期窗口；系统把泛泛的监管截止日判为空泛锚点，除非来源明确给出公司的预约披露日期。

锚点的作用域分为公司、行业、市场和监管四类。作用域的作用是防止跨作用域误用。例如，同行业其他公司的营收指引结果可以作为行业锚点，但不能直接当作目标公司的结果；交易所的披露截止日可以作为监管锚点，但不能直接推断公司会在截止日当天披露。资格判定检查作用域是否与断言匹配，不匹配的锚点被判为对抗锚点。

\section{与既有框架的区别}

FinPlanRAG与断言验证、RAG和选择性预测的区别可以概括为三点。

第一，状态不同。断言验证的输出是一个离散裁决，RAG的输出是一个生成答案，选择性预测的输出是一个标签或弃权。FinPlanRAG的输出是一个可演化的锚点状态，包含证据、资格、可靠性和信念。这个状态可以被逐条检查，也可以在结果揭示后更新。

第二，时间不同。断言验证和RAG通常在一次调用中完成，不区分证据的可用时间和结果的揭示时间。FinPlanRAG把截止时点写进契约，把揭示日写进可靠性更新规则。这个区分使系统可以在回测中复现部署时的信息条件。

第三，决策不同。断言验证和RAG通常报告准确率或F1，选择性预测报告覆盖率或弃权率。FinPlanRAG在同一个损失下比较点答案、裁决集和弃权，报告决策损失、覆盖率和集合大小。这个口径使系统可以在不确定性高时选择集合或弃权，而不是被迫输出一个点答案。

% ==================== 第5章 ====================

本章给出FinPlanRAG中AnchorEvolve模块的完整算法设计。本章第1节到第7节分别介绍证据构造、锚点抽取、资格判定、时间演化、信念更新、全动作读出和披露日期断言，最后给出算法伪代码、复杂度和审计格式。

\section{证据构造}

证据构造的目标是为每条断言提供一组带日期、带来源编号、可引用的材料。构造过程分为查询生成、检索、去重、日期过滤和编号五步。

查询生成从契约出发。给定公司$e$、目标期间$p$和指标$m$，系统生成四类查询：公司加期间加指标，公司加最近季度加指标，公司加上一年度加指标，以及公司加“预测”“研报”“公告”等限定词。对于披露日期断言，系统额外生成公司加“年度报告披露日期”、公司加“往年披露时间”、公司加“预约披露”等查询。查询数量由配置限制，默认每条断言最多四条。查询中不放入断言的目标答案，避免把判断泄漏到检索。

检索使用多个网页搜索提供者，并按轮换顺序调用。某个提供者返回空结果或安全验证页时，系统自动切换到下一个提供者。这个设计是为了避免单一搜索引擎的限流让整次运行静默失败。检索结果包含标题、链接、摘要、发布日期和提供者。对于披露日期断言，系统还查询巨潮资讯网的公告接口，直接获取公司历年年报的公告日期和PDF链接。官方公告记录作为高优先级证据，排在网页结果之前。

去重按链接和标题进行。同一篇文章的不同转载只保留一条；同一链接的不同摘要只保留一条。日期过滤检查证据的发布日期是否晚于截止时点。晚于截止时点的证据被写入剔除列表，不进入锚点抽取。发布日期未知的证据保留，但标记为无日期，并在资格判定中降权。

编号使用$W1,W2,\dots$，按证据得分排序。得分由日期是否已知、摘要是否包含数字、标题或摘要是否包含公司名、文本长度和查询顺序决定。官方公告记录得分最高。编号的作用是让锚点可以引用具体来源，也让审计报告可以把锚点映射回证据。

对于申报文件，系统在证据构造阶段还做段落切分。每个段落带稳定编号$P1,P2,\dots$，结构化证据带$E1,E2,\dots$。抽取器可以引用$P\#$或$E\#$，资格判定检查引用的段落是否在时点摘录中。这个设计使锚点可以定位到段落，而不是只定位到文件。

\section{锚点抽取}

锚点抽取是一个LLM调用。输入包括断言契约、结果标签及其定义、参考口径、时点证据、锚点类型和关系类型。输出是一个JSON对象，包含锚点列表、标签分数、日期估计（仅日期断言）、置信度和缺失说明。

提示词要求锚点遵守以下规则。第一，每条锚点必须引用一个证据编号，不能引用不存在的来源。第二，每条锚点必须给出逐字引文，引文必须在来源文本中出现。第三，锚点的发布日期不得晚于截止时点；来源日期未知时，日期设为空，可靠性不得高于0.55。第四，可靠性是“隐含标签正确的校准概率”，不是推理置信度。第五，硬锚点必须是显式约束，能够排除至少一个标签；管理层评论和分析师预测通常是软锚点。第六，证据不足时在缺失字段中说明，不编造事实。第七，标签分数必须使用契约中的标签，且和为1。

抽取器的输出是一个结构化对象，而不是自由文本。这个约束使后续的资格判定可以按字段检查，而不是依赖另一个LLM来判断前一个LLM的输出是否合理。解析失败的响应被记录为错误，不进入信念更新；系统不使用金标结果作为回退。这个设计是为了避免在评估中把解析失败悄悄地替换成正确答案。

锚点抽取的成本主要来自输入和输出长度。单条断言的输入通常在数千token量级，输出在数百token量级。系统按请求负载的哈希缓存响应，相同输入产生相同输出。缓存不包含API密钥，也不写入审计文件的公开部分。

\section{五类资格判定}

资格判定是FinPlanRAG与普通RAG的关键区别。普通RAG把检索到的片段直接交给生成器；FinPlanRAG在锚点进入信念更新之前，先用确定性规则检查锚点的资格。表\ref{tab:qualification}给出五类判定。

\begin{table}[H]
\centering
\zihao{5}
\caption{五类锚点资格判定。}
\label{tab:qualification}
\begin{tabular}{p{2.1cm}p{7.0cm}p{4.2cm}}
\toprule
类型 & 判定条件 & 处理 \\
\midrule
空泛 & 没有隐含标签，或锚点不提供有效标签信息 & 不进入更新 \\
过期 & 锚点日期或来源日期晚于截止时点 & 不进入更新 \\
对抗 & 来源编号不在证据集合中，或作用域不匹配 & 不进入更新 \\
硬 & 显式约束、引文通过、来源日期已知、可靠性不低于阈值 & 约束可行集 \\
软 & 其余有效锚点 & 似然更新 \\
\bottomrule
\end{tabular}
\end{table}

资格判定的具体规则如下。

第一，来源检查。锚点的来源编号必须出现在本次证据集合中。来源不存在，判为对抗。这个检查防止LLM引用一个看起来合理但实际不存在的来源。

第二，日期检查。锚点日期或来源日期晚于截止时点，判为过期。如果锚点日期与来源日期不一致，系统以来源日期为准检查过期，并对可靠性乘以0.9。日期不一致本身不是拒绝理由，但它是一个时间一致性疑点。

第三，引文检查。锚点的引文必须在来源文本中逐字出现。系统对引文和来源文本做空白规范化，然后做子串匹配。引文为空、长度过短或匹配失败，判为引文未通过。未通过的锚点可靠性乘以0.7，且不能成为硬锚点。这个检查是确定性的，不依赖另一个LLM。

第四，作用域检查。锚点的作用域必须是公司、行业、市场或监管之一。作用域不匹配，判为对抗。例如，把同行业其他公司的结果当作目标公司的直接结果，作用域不匹配。

第五，无日期检查。锚点日期和来源日期都未知时，锚点保留为软锚点，但可靠性上限为0.55，且不能成为硬锚点。这个规则承认无日期证据可能有用，但不允许它主导判断。

第六，参考值检查。对于增速类断言，只描述上一年度参考值的锚点判为空泛。上一年度实际值是契约的参考口径，不是目标期间的结果证据。这个检查防止抽取器把基数当作结果，从而把信念拉向低增长。

第七，监管规则检查。对于披露日期断言，泛泛的法定截止日判为空泛。截止日是一个上界，不能区分三个日期窗口；只有公司预约披露日期或历史披露记录才能支持具体窗口。

第八，日期标签投影。对于披露日期断言，如果锚点提到的日期与隐含标签不一致，系统把日期中的月和日投影到目标年份，再匹配日期窗口。例如，历史年报日期2025-04-21投影到2027-04-21，落在“晚于2027-04-15”窗口。投影只在一个锚点提到的所有日期指向同一窗口时生效；如果锚点同时提到多个窗口的日期，系统不覆盖模型标签，只记录不一致。这个规则修正了抽取器在跨年日期上的常见错误。

第九，时间一致性检查。对于披露日期断言，如果引文中提到的日期晚于来源发布日期，且不晚于截止时点，系统标记时间一致性疑点，可靠性乘以0.5，且不能成为硬锚点。这个检查处理动态网页或转载文章在原始发布日期之后被更新，导致引文包含当时不可能知道的信息的情况。

资格判定的输出是每条锚点的状态、资格类型、有效可靠性、引文是否通过、时间一致性标记和理由。被拒绝的锚点保留在审计文件中，理由逐条列出；被采纳的锚点进入信念更新。这个设计使审计者可以看到系统拒绝了什么，以及为什么拒绝。

\section{时间演化}

时间演化的目标是让锚点类型的可靠性随已揭示结果更新，同时不把未来信息泄漏到时点判断。设断言$q$的截止时点为$c$，锚点类型为$\tau$。截至$c$已揭示的、类型为$\tau$的锚点中，隐含标签与结果一致的数量为$m_\tau$，总数为$n_\tau$。类型$\tau$的Beta可靠性为
\begin{equation}
r_\tau=\frac{m_\tau+\alpha}{n_\tau+\alpha+\beta},
\label{eq:beta}
\end{equation}
其中$\alpha,\beta$是Beta先验参数。默认取$\alpha=\beta=1$，即均匀先验。当$n_\tau=0$时，$r_\tau=0.5$。

抽取器给出的可靠性是LLM对单条锚点的主观判断，不是经过验证的频率。系统把演化可靠性与LLM可靠性做收缩：
\begin{equation}
\rho_\tau=\lambda r_\tau+(1-\lambda)r_{\mathrm{LLM}},
\label{eq:shrink}
\end{equation}
其中$\lambda$是收缩系数，默认0.5。$\rho_\tau$作为软锚点更新的指数。这个收缩的作用是：当历史样本少时，系统更多依赖LLM判断；当历史样本多时，系统更多依赖已揭示结果的频率。

先验也按揭示日更新。设截至$c$已揭示的结果中，落在标签$h$的数量为$n_h$，总数为$n$，则先验为
\begin{equation}
b_0(h)=\frac{n_h+1}{n+|\mathcal H|}.
\label{eq:prior}
\end{equation}
当$n=0$时，先验为均匀分布。这个先验只使用已揭示结果，不使用目标期间的结果。

演化状态按轮次更新。第$r$轮的结果在第$r$轮全部断言裁决完成后才更新状态，并从第$r+1$轮开始使用。新发现的锚点类型的生效轮次为$r+1$，因此触发发现的断言不能使用该类型。这个规则保证类型库扩展是时点正确的。本文的主实验报告单轮扩展；多轮扩展作为接口保留，第\ref{ch:conclusion}章讨论其局限。

\section{类型库的验证集扩展}

固定锚点类型库依赖人工设计，难以覆盖不同指标和不同市场的证据模式。本文在Beta可靠度演化之外加入一个受验证集约束的类型库扩展步骤。它不是让模型自由改写标签或测试集，而是让模型在验证误差上提出新的锚点类型，只有通过验证集损失检验的类型才进入类型库。

设第$t$轮的锚点类型库为$\mathcal T_t$，验证集上由当前类型库得到的决策损失为$L_{\mathrm{val}}(\mathcal T_t)$。循环分四步。

第一步，误差挖掘。把验证集按公司不重叠地分成挖掘集和选择集。在挖掘集上，用当前类型库计算每条断言的预测和损失，选出预测错误或损失较高的断言。挖掘集只用于提出候选类型，不参与候选接受。

第二步，类型提出。把错误断言、对应证据、当前锚点和真实结果交给一个LLM，要求它提出至多$k$个新的锚点类型。每个候选类型是一个四元组
\begin{equation}
\tau'=(\text{名称},\ \text{L0关系},\ \text{字段},\ \text{适用条件}),
\label{eq:candidatetype}
\end{equation}
并附带一句证据示例。提出器被要求只描述可复用的证据模式，不能引用具体测试断言，不能把结果标签本身当作类型。

第三步，候选重抽取。把候选类型加入类型库，重写锚点抽取提示词中的类型说明，在选择集上重新抽取锚点并计算后验。类型库扩展为$\mathcal T_t\cup\mathcal T'$。

第四步，验证门控。只有当选择集上的损失下降超过阈值$\delta$时，才接受候选类型：
\begin{equation}
L_{\mathrm{val}}(\mathcal T_t)-L_{\mathrm{val}}(\mathcal T_t\cup\mathcal T')>\delta.
\label{eq:typeaccept}
\end{equation}
否则拒绝整组候选，保持原类型库。接受后进入下一轮，最多进行$R$轮。测试集只在循环冻结后评估一次，不参与类型提出、阈值选择或轮数选择。

这个循环与提示词演化分开评估。提示词改写、类型扩展和二者组合分别作为候选臂，在同一选择集上比较；只有损失最低且超过阈值的臂被接受。这样可以把“模型提出了新类型”和“新类型真的改善了验证损失”分开，避免把语言上的合理性当作经验收益。

类型库扩展的代价是额外的验证集抽取。每轮需要重抽选择集锚点，但不重抽测试集；类型库冻结后才抽取测试集。由于循环只使用验证集，它与主实验的时点约束一致，也不会把测试结果泄漏到类型库。

类型库扩展是框架的训练步骤，不是独立模块。算法\ref{alg:typeexpand}给出伪代码，其输出$\mathcal T^\star$作为算法\ref{alg:anchorevolve}的类型库输入。主实验中DeepSeek使用通过验证门控的类型库，GLM和Qwen保留基础库；这样，类型库的选择和其他超参数一样，只由验证集决定。

\begin{algorithm}[H]
\caption{类型库的验证集扩展}
\label{alg:typeexpand}
\begin{algorithmic}[1]
\Require 类型库$\mathcal T_0$，验证集$\mathcal D_{\mathrm{val}}$，阈值$\delta$
\Ensure 冻结类型库$\mathcal T^\star$
\State 将$\mathcal D_{\mathrm{val}}$按公司不重叠划分为挖掘集$\mathcal D_{\mathrm{mine}}$和选择集$\mathcal D_{\mathrm{select}}$
\State 在$\mathcal D_{\mathrm{mine}}$上计算当前类型库的损失，选出高损失断言
\State 由LLM提出候选类型$\mathcal T'$，每个类型包含关系、字段和适用条件
\State 在$\mathcal D_{\mathrm{select}}$上重抽锚点，计算$L_{\mathrm{val}}(\mathcal T_0\cup\mathcal T')$
\If{$L_{\mathrm{val}}(\mathcal T_0)-L_{\mathrm{val}}(\mathcal T_0\cup\mathcal T')>\delta$}
\State $\mathcal T^\star\gets\mathcal T_0\cup\mathcal T'$
\Else
\State $\mathcal T^\star\gets\mathcal T_0$
\EndIf
\State \Return $\mathcal T^\star$ \Comment{冻结后用于测试集}
\end{algorithmic}
\end{algorithm}


\section{信念更新}

信念从先验$b_0$开始。对于软锚点$a$，设其隐含标签为$\hat h$，有效可靠性为$\rho$，软似然为
\begin{equation}
L_a(h)=
\begin{cases}
p, & h=\hat h,\\
\dfrac{1-p}{|\mathcal H|-1}, & h\ne\hat h,
\end{cases}
\label{eq:softlikelihood}
\end{equation}
其中$p$是配置项，默认0.7。软锚点的更新为
\begin{equation}
b(h)\propto b(h)\,L_a(h)^{\rho}.
\label{eq:softupdate}
\end{equation}
指数$\rho$控制锚点的影响强度：可靠性越高，更新越强；可靠性越低，更新越弱。更新后对信念做归一化。

对于硬锚点，系统把信念压到隐含标签上，并保留一个小底：
\begin{equation}
b(h)=
\begin{cases}
1-\epsilon, & h=\hat h,\\
\dfrac{\epsilon}{|\mathcal H|-1}, & h\ne\hat h,
\end{cases}
\label{eq:hardupdate}
\end{equation}
其中$\epsilon$默认0.02。小底的作用是避免后续软锚点无法恢复。多个硬锚点的可行集交集、冲突检测和回退到软更新是论文的设计目标，当前实现只处理单硬锚点的点质量更新；第\ref{ch:conclusion}章把这一点列为局限。

锚点后验记为$b_A$。抽取器还输出标签分数$s$，它是边际信号，不是锚点状态的替代。系统把$b_A$与$s$按权重$\gamma$混合：
\begin{equation}
b_A'=\gamma b_A+(1-\gamma)s,
\label{eq:blend}
\end{equation}
其中$\gamma$默认0.5。混合后的后验保留锚点的结构化信息，同时利用抽取器对三类结果的边际判断。

结构化随机森林后验记为$b_R$。输入特征包括指引宽度、当前季度值相对指引中点的比例、指引与当前季度的比值、历史结果计数与频率、前次指引增长与宽度、目标季度、截止年份、行业，以及演化先验特征（同行业历史、公司意外历史、日历、当前季度相对上下界的位置、可靠性标志）。随机森林在开发集与验证集上训练，并在公司不重叠折上做温度缩放；测试公司不出现在训练集中。

锚点后验与随机森林后验融合：
\begin{equation}
b_F=w b_A'+(1-w)b_R,
\label{eq:fusion}
\end{equation}
其中$w$只在验证集上从配置的权重网格中选择，测试集不参与选择。默认网格为$\{0,0.1,0.2,0.3\}$。最终实验中选择的权重在正文中报告。

\section{全动作损失感知读出}

给定融合后验$b_F$，系统在动作空间$\mathcal A$上选择期望损失最小的动作。动作空间包括$|\mathcal H|$个点答案、$\sum_{k=2}^{|\mathcal H|}\binom{|\mathcal H|}{k}$个非单元素裁决集，以及弃权。对于三标签结果空间，动作空间包含3个点答案、4个非单元素裁决集和1个弃权，共8个动作。

点答案的期望损失为
\begin{equation}
\E[L_{\text{ans}}(h)]=\sum_{y\in\mathcal H}b_F(y)L_{\text{ans}}(h,y).
\label{eq:eans}
\end{equation}
裁决集的期望损失为
\begin{equation}
\E[L_{\text{set}}(R)]=\sum_{y\in\mathcal H}b_F(y)L_{\text{set}}(R,y).
\label{eq:eset}
\end{equation}
弃权的期望损失为$\ell_{\text{abstain}}$。系统选择
\begin{equation}
a^\star=\arg\min_{a\in\mathcal A}\E[L(a)].
\label{eq:fullaction}
\end{equation}

全动作读出与受限可信集门的区别在于动作空间。受限门只比较点答案和最高密度$1-\alpha$可信集，可信集的大小由覆盖率目标决定，而不是由损失决定。全动作读出把点答案、裁决集和弃权放在同一个损失下比较，因此可以在不确定性高时选择裁决集或弃权，也可以在置信度高时选择点答案。对于同一个后验和同一个损失，全动作读出的期望损失不会高于受限门，因为受限门的动作是全动作空间的子集。这个性质是H3的形式基础。

单元素裁决集不进入动作空间。原因是单元素裁决集与点答案在语义上重复，但在损失上不同：如果漏报惩罚低于错误分类惩罚，单元素裁决集会系统性地优于点答案，导致系统把点答案都标成裁决集。这个现象不是决策改进，而是动作标签的重复计数。本文把裁决集限制为两个及以上标签，使三类动作的含义保持清晰。

可信集$R_\alpha$仍然作为覆盖口径报告：
\begin{equation}
R_\alpha=\arg\min_{R\subseteq\mathcal H}\{|R|:\sum_{h\in R}b_F(h)\ge1-\alpha\}.
\label{eq:credible}
\end{equation}
它不参与动作选择，但用于审计和覆盖率报告。损失最优裁决集$R^\star$是动作选择的结果，可以在$R_\alpha$更宽时选择更窄的集合，也可以在不确定性高时选择弃权。两者在审计报告中分开列出，避免把覆盖率误读为决策质量。

\section{披露日期断言}

披露日期断言的算法与数值指标断言共享同一套锚点状态和读出，但结果空间和日期估计不同。

结果空间由日历窗口定义。给定目标期间$p$，系统根据交易所规则和公司历史披露习惯构造三个窗口。例如，2026年年度报告的窗口为“不晚于2027-03-31”“2027-04-01至2027-04-15”“晚于2027-04-15”。窗口必须互斥且覆盖法定披露截止日之前的全部日期。

锚点抽取器除了输出锚点和标签分数，还输出一个日期估计，包括点估计、区间下界、区间上界和依据。日期估计是辅助信息，不直接决定动作；动作仍由标签后验和损失读出决定。系统把日期估计与标签后验做一致性检查：如果点估计落在最高信念窗口内，系统保留点估计；如果点估计落在窗口外，系统用最高信念窗口的中点作为回退点估计。区间下界和上界优先使用模型给出的值，但限制在合理范围内。

日期标签投影是披露日期断言的关键修正。历史年报日期、预约披露日期和监管截止日往往跨年，抽取器容易把历史日期直接映射到当前窗口。系统把日期中的月和日投影到目标年份，再匹配窗口。投影规则只在一个锚点提到的所有日期指向同一窗口时生效，避免把包含多个日期的锚点强行归类。

披露日期断言的评价不以点日期误差作为主指标，因为公告日期与实际披露时点可能相差一天。本文在A股回测上报告日期窗口准确率，并在案例研究中展示动作、裁决集、点估计、区间和审计轨迹。

\section{算法与复杂度}

算法\ref{alg:anchorevolve}给出单条断言的完整流程。

\begin{algorithm}[H]
\caption{AnchorEvolve单条断言的证据构造、锚点演化与全动作裁决}
\label{alg:anchorevolve}
\begin{algorithmic}[1]
\Require 断言契约$q=(e,p,m,\mathcal H,c)$，演化状态$S_c$，类型库$\mathcal T^\star$，随机森林后验$b_R$
\Ensure 动作$a^\star$，点答案$h^\star$，裁决集$R^\star$，后验$b_F$
\State $V_c\gets\textsc{CollectEvidence}(q,c)$ \Comment{按截止时点过滤，分配$W\#$}
\State $A\gets\textsc{ExtractAnchors}(q,V_c,S_c,\mathcal T^\star)$ \Comment{LLM，要求来源与逐字引文}
\State $A_q\gets\textsc{Qualify}(A,q,c)$ \Comment{空泛/过期/对抗/硬/软}
\State $b_A\gets\textsc{BeliefUpdate}(S_c,A_q)$
\State $b_A'\gets\gamma b_A+(1-\gamma)\textsc{LabelScores}(A)$
\State $b_F\gets w b_A'+(1-w)b_R$
\State $(a^\star,h^\star,R^\star)\gets\textsc{FullActionReadout}(b_F)$
\State $S_{c^+}\gets\textsc{RevealUpdate}(S_c,c)$ \Comment{仅用揭示日早于下一截止时点的结果}
\State \Return $(a^\star,h^\star,R^\star,b_F)$
\end{algorithmic}
\end{algorithm}

算法的时间成本由三部分组成。第一是证据构造，包括网页检索、页面抓取和官方公告查询。网页检索的数量由配置限制，默认每条断言四条查询；页面抓取只抓前若干条结果，避免不必要的网络开销。第二是锚点抽取，一次LLM调用的输入在数千token量级，输出在数百token量级。第三是读出，动作空间的大小为$|\mathcal H|+\sum_{k=2}^{|\mathcal H|}\binom{|\mathcal H|}{k}+1$，当$|\mathcal H|\le6$时最多63个动作，计算可以忽略。随机森林推断在本地完成，不产生API成本。

系统的存储成本主要来自审计文件。每条断言保存证据列表、锚点列表、信念轨迹、动作损失和LLM响应摘要。审计文件的大小通常在数百KB量级，适合长期保存和人工复核。LLM响应缓存单独存放，不进入公开审计文件。

\section{可审计性}

可审计性是FinPlanRAG的设计目标之一。系统把每个中间状态写入结构化文件，使读者可以从最终动作回溯到证据和锚点。表\ref{tab:auditfields}给出审计文件的主要字段。

\begin{table}[H]
\centering
\zihao{5}
\caption{审计文件的主要字段。}
\label{tab:auditfields}
\begin{tabular}{p{2.8cm}p{8.8cm}}
\toprule
对象 & 字段 \\
\midrule
契约 & 公司、代码、目标期间、指标、结果空间、标签定义、参考口径、截止时点 \\
证据 & 来源编号、标题、链接、发布日期、提供者、摘要、正文、是否官方来源 \\
锚点 & 类型、取值、来源编号、逐字引文、发布日期、可靠性、资格、隐含标签、理由、纠正记录 \\
信念 & 先验、仅锚点后验、融合后验、逐步更新轨迹 \\
动作 & 动作类型、点答案、裁决集、可信集、期望损失、各动作损失排序 \\
日期 & 点估计、区间、最高信念窗口、模型估计与窗口是否一致、依据 \\
\bottomrule
\end{tabular}
\end{table}

审计文件有两个用途。第一，复核。分析师可以逐条检查被采纳的锚点是否真的支持隐含标签，被拒绝的锚点是否应当被拒绝，证据的发布日期是否早于截止时点。第二，问责。当判断出现错误时，系统可以定位错误来自证据缺失、锚点抽取、资格判定、信念更新还是动作选择，而不是把错误归因于一个黑箱模型。

可审计性不等于正确性。一个锚点即使有来源和引文，也可能是对来源的误读；一个资格判定即使规则明确，也可能过严或过松。本文在案例研究中展示审计轨迹，并在第\ref{ch:conclusion}章把人工审计实验列为未来工作。当前的可审计性是设计属性，不是已经测量的用户收益。

